\documentclass[11pt]{article}

\usepackage[margin=0.75in]{geometry}
\usepackage[T1]{fontenc}
\usepackage[utf8]{inputenc}
\usepackage{enumitem}
\usepackage[hidelinks]{hyperref}
\usepackage{titlesec}

\pagestyle{empty}
\setlength{\parindent}{0pt}
\setlength{\parskip}{0pt}
\titleformat{\section}{\large\bfseries}{}{0pt}{}[\titlerule]
\titlespacing*{\section}{0pt}{0.8em}{0.4em}
\setlist[itemize]{leftmargin=*, noitemsep, topsep=0pt}

\begin{document}

\begin{center}
	{\LARGE \textbf{Kaden Downes}}\\
	Fort Wayne, IN  $\vert$  (260) 760-9476 $\vert$ \href{mailto:zaetep@hotmail.com}{zaetep@hotmail.com} $\vert$ \href{https://github.com/zaetep}{github.com/zaetep}
\end{center}

\section*{Summary}
Aspiring Embedded Software Engineer with hands-on experience in ARM Cortex-M firmware development, FreeRTOS, PCB design, and hardware debugging. Strong background in low-level computer hardware, communication protocols, and cross disciplinary engineering projects. Available to begin full-time July 2026. Authorized to work in the U.S. without sponsorship. Willing to relocate as needed.

\section*{Skills}
\begin{itemize}
	\item Embedded Systems
	\item FreeRTOS and related concepts
	\item PCB and RTL Design
	\item Communication Protocols (UART, I2C, SPI, Wi-Fi, Bluetooth)
	\item Hardware Debugging (J-Link, Logic Analyzers)
	\item \textbf{Tools:} Altium, MATLAB, Git, Azure DevOps, Vivado, Microsoft Office
	\item \textbf{Embedded Languages:} C/C++, Rust
	\item \textbf{Other Languages:} C\#, Python, VHDL, Verilog, Dart, LaTeX
	\item \textbf{Other Skills:} Linux, Networking, Troubleshooting, System Administration
\end{itemize}

\section*{Education}
\begin{tabular*}{\textwidth}{@{\extracolsep{\fill}} l c r}
	\textbf{Bachelor of Science in Electrical Engineering} & Purdue University Fort Wayne & 2022--2026\\
	\textbf{Bachelor of Science in Computer Engineering} & Purdue University Fort Wayne & 2022--2026
\end{tabular*}

% \section*{Relevant Coursework}
% \begin{itemize}
% 	\item Embedded Microprocessors --- ARM architecture, peripheral interfacing, low-level firmware development
% 	\item Real-Time Operating Systems (FreeRTOS) --- task scheduling, ISR design, synchronization primitives
% 	\item Computer Architecture --- pipelining, memory hierarchy, and performance optimization for constrained systems
% \end{itemize}

\section*{Projects}
\textbf{Battery Powered Wet Floor Sensor (Senior Design)} \hfill \textit{Sponsor: Franklin Electric}\\
	Implemented device firmware in C for an ARM Cortex-M MCU using FreeRTOS; managed tasks, ISRs, and power states.
	Integrated Wi-Fi module with Bluetooth provisioning and implemented TCP communication for telemetry and alerts.
	Built project with CMake/Ninja and managed source with Git. Designed PCB in Altium for manufacturing. Supported hardware bring-up and debugging with J-Link and logic analyzers.\\[0.25em]
\textbf{Blinky (Rust Proof-of-Concept)} \hfill \href{https://github.com/zaetep/blinky}{GitHub Link}\\
	A simple proof-of-concept program that blinks the LED on the NUCLEO-U545RE-Q STM32 development board. This project is written entirely in Rust, and uses the Embassy HAL and executor to blink an LED. The project is set up to use cargo embed, which automates the build and flash process.\\[0.25em]
\textbf{IoT Sensor (In Progress)} \hfill \href{https://github.com/zaetep/iot_sensor}{Firmware} | \href{https://github.com/zaetep/python_iot_receiver_gui}{Receiver}\\
	Applying my previous experience, this project consists of a firmware portion and a receiver portion. The firmware is written on STM32, where it will collect data from sensors and transmit the data over Wi-Fi via an ESP32-C6. The firmware is written in C using FreeRTOS, with CMake as the build system. The receiver is a Python GUI program that will receive the packets sent by the firmware and display the data.

% \section*{Certifications}
% \begin{itemize}
% 	\item Google IT Support Professional Certificate (v2) --- Coursera, 2025
% \end{itemize}

% \textbf{Project Name} \hfill \href{https://github.com/yourrepo}{GitHub/Link}\\
% Brief description of the project, technologies used, and key impact or outcome.

\section*{Experience}
\textbf{Advanced Repair Agent} \hfill Geek Squad --- Fort Wayne \hfill Nov 2021--Present
\begin{itemize}
	\item Performed detailed fault analysis to determine software vs. hardware causes and carried out component level repairs on multiple device platforms.
	\item Developed automation scripts using PowerShell and Bash to streamline workflows.
	\item Closed 137 tickets every month, averaging 1.2 tickets per hour with 140\% labor efficiency.
	\item Communicated with clients and other agents as to the status of and next steps for work orders.
	\item Responsibly handled sensitive customer data, including storing and transferring personal files.
	\item Achieved 100\% customer satisfaction according to Apple repair surveys.
\end{itemize}



\end{document}
